% УВОД, без номер.
\section*{УВОД}
\label{sec:introduction}
\addcontentsline{toc}{section}{УВОД}

Достъпът до структурирани данни изисква владеене на формален език за заявки. Човек,
който знае какво търси в една база от данни, често не знае как да го напише на SQL.
Интерфейсът на естествен език премахва точно тази бариера.

Настоящата курсова работа разработва такъв интерфейс под името \term{Lexicon}. Системата
приема въпрос на естествен език, анализира го морфологично, синтактично и семантично, и
генерира заявка на SQL или SPARQL срещу зададена предметна област. Обработката е изцяло
символна: всяка стъпка от входното изречение до изходната заявка е представима, проверима
и обяснима.

\textbf{Целта} е работеща верига от анализ и синтез, която превръща въпрос във формална
заявка и връща на потребителя перифраза на разбрания въпрос, за да може той да потвърди
или отхвърли тълкуването, преди заявката да бъде изпълнена. Оттук произтичат задачите:
морфологичен анализатор върху краен преобразувател, синтактичен анализатор, който запазва
всички разчитания, семантично изображение върху модел на знанието, генератор на формална
заявка, синтез на потвърждаваща перифраза и методика за експериментална проверка.

\textbf{Обхватът} е ограничен до въпросителни изречения от предметната област, описана в
лексикона и в модела на знанието. Извън обхвата остават свободен диалог, анафора през
няколко реплики и произволна предметна област без предварително описание.

Раздел~\ref{sec:theory} излага теоретичната основа и обосновава избора за всеки слой.
Раздел~\ref{sec:design} представя проектирането, Раздел~\ref{sec:implementation}
програмната реализация, а Раздел~\ref{sec:results} методиката за измерване и получените
резултати.
